% 第2章 2.1.3：标准正交基前 n 项投影 P_n 的强收敛与范数不收敛。
% x 的系数 c_k=2^{-k/2}（‖x‖=1），尾误差 ‖x-P_nx‖=2^{-n/2}；
% y 的系数 c_k=(sqrt6/π)/k（‖y‖=1），尾误差 (sqrt6/π)(Σ_{k>n}1/k^2)^{1/2}，由 Python 数值求和；
% 对被截掉的单位向量 e_{n+1}，‖(I-P_n)e_{n+1}‖=1=‖I-P_n‖。
\begin{tikzpicture}[font=\figurefont\footnotesize]
\pgfplotsset{every axis/.style={width=54mm,height=38mm,axis lines=left,
  x axis line style={draw=FigureStroke,line width=.6pt},y axis line style={draw=FigureStroke,line width=.6pt},
  tick style={draw=FigureStroke},tick label style={font=\figurefont\scriptsize,text=FigureInk},
  label style={font=\figurefont\scriptsize},
  xlabel style={at={(axis description cs:1,0)},anchor=north west,yshift=1mm},
  ymajorgrids,grid style={draw=FigureGrid,line width=.4pt},
  title style={font=\figurefont\footnotesize},
  legend style={draw=none,fill=none,font=\figurefont\scriptsize,cells={anchor=west}}}}
\begin{axis}[name=A,ybar,bar width=1.8mm,xlabel={$k$},xmin=0.3,xmax=10.7,ymin=0,ymax=0.8,
  xtick={1,2,...,10},ytick={0,0.25,0.5,0.75},title={(a) $x$ 的系数，$n=4$}]
\addplot[fill=FigureBlue,draw=FigureBlue] coordinates {(1,0.7071) (2,0.5) (3,0.3536) (4,0.25)};
\addplot[fill=FigureBlueFill,draw=FigureBlue] coordinates {(5,0.1768) (6,0.125) (7,0.0884) (8,0.0625) (9,0.0442) (10,0.0312)};
\node[font=\FigureNoteFont\fontsize{7.5}{9}\selectfont,text=FigureMuted,anchor=south] at (axis cs:3.5,0.40) {保留};
\node[font=\FigureNoteFont\fontsize{7.5}{9}\selectfont,text=FigureMuted,anchor=south] at (axis cs:7.5,0.25) {尾项 $x-P_4x$};
\end{axis}
\begin{axis}[at={(A.east)},anchor=west,xshift=10mm,xlabel={$n$},xmin=0,xmax=30,ymin=0,ymax=1.15,
  xtick={0,10,20,30},ytick={0,0.5,1},title={(b) 截断误差},
  legend style={at={(0.5,-0.3)},anchor=north},legend columns=3]
\addplot[draw=FigureBlue,line width=1.1pt,mark=*,mark size=.8pt] coordinates {(0,1.0000) (1,0.7071) (2,0.5000) (3,0.3536) (4,0.2500) (5,0.1768) (6,0.1250) (7,0.0884) (8,0.0625) (9,0.0442) (10,0.0312) (11,0.0221) (12,0.0156) (13,0.0110) (14,0.0078) (15,0.0055) (16,0.0039) (17,0.0028) (18,0.0020) (19,0.0014) (20,0.0010) (21,0.0007) (22,0.0005) (23,0.0003) (24,0.0002) (25,0.0002) (26,0.0001) (27,0.0001) (28,0.0001) (29,0.0000) (30,0.0000)};
\addplot[draw=FigureGreen,line width=1.1pt,dashed] coordinates {(0,1.0000) (1,0.6262) (2,0.4900) (3,0.4154) (4,0.3668) (5,0.3320) (6,0.3055) (7,0.2845) (8,0.2673) (9,0.2529) (10,0.2405) (11,0.2298) (12,0.2205) (13,0.2122) (14,0.2047) (15,0.1980) (16,0.1919) (17,0.1864) (18,0.1813) (19,0.1765) (20,0.1722) (21,0.1681) (22,0.1644) (23,0.1608) (24,0.1575) (25,0.1544) (26,0.1515) (27,0.1487) (28,0.1460) (29,0.1435) (30,0.1412)};
\addplot[draw=FigureYellow,line width=1.1pt,dash dot] coordinates {(0,1) (30,1)};
\legend{$\|x-P_nx\|$,$\|y-P_ny\|$,$\|(I-P_n)e_{n+1}\|$}
\end{axis}
\end{tikzpicture}
